\ifdefined\gkver\else\def\gkver{student}\fi
\documentclass[\gkver]{gaokaozhenti}
\begin{document}

\chapter{2011年天津理综}
%% paperType: 理综物理部分

\begin{enumerate}
\item
%% number: 1
%% paperName: 2011年天津理综第 1 题 6 分
%% typeId: 1
%% type: 单选题
%% chapter: 原子和原子核
%% point: 原子核式结构
%% method: 无
%% score: 6
%% degree: 650
%% duplicateId: 0
%% body:
下列能揭示原子具有核式结构的实验是\xzanswer{C}
\fourchoices[answer=C]
{光电效应实验}
{伦琴射线的发现}
{$\alpha$粒子散射实验}
{氢原子光谱的发现}
%% answer: C
%% memo:
\memoanswer{
光电效应实验说明光具有粒子性，A 错误；X 射线（伦琴射线）的发现是 19 世纪末 20 世纪初物理学的三大发现（X 射线 1896 年、放射线 1896 年、电子 1897 年）之一，这一发现标志着现代物理学的产生，B 错误；氢原子光谱的发现解释了原子的稳定性以及原子光谱的分立特征，D 错误；卢瑟福的 $\alpha$ 粒子散射实验揭示原子具有核式结构，C 正确。
}

\item
%% number: 2
%% paperName: 2011年天津理综第 2 题 6 分
%% typeId: 1
%% type: 单选题
%% chapter: 运动和力
%% point: 连接体
%% method: 无
%% score: 6
%% degree: 650
%% duplicateId: 0
%% body:
如图所示，A、B 两物块叠放在一起，在粗糙的水平面上保持相对静止地向右做匀减速直线运动，运动过程中 B 受到的摩擦力\xzanswer{A}

\onepicture[w=4.5cm]{figs/02.png}

\fourchoices[answer=A]
{方向向左，大小不变}
{方向向左，逐渐减小}
{方向向右，大小不变}
{方向向右，逐渐减小}
%% answer: A
%% memo:
\memoanswer{
对于多个物体组成的物体系统，若系统内各个物体具有相同的运动状态，应优先选取整体法分析，再采用隔离法求解。取 A、B 系统整体分析有 $f_{\text{地}A}=\mu(m_A+m_B)g=(m_A+m_B)a$，$a=\mu g$，B 与 A 具有共同的运动状态，取 B 为研究对象，由牛顿第二定律有：$f_{AB}=\mu m_B g=m_B a=$ 常数，物体 B 做速度方向向右的匀减速运动，故而加速度方向向左。A 正确。
}

\item
%% number: 3
%% paperName: 2011年天津理综第 3 题 6 分
%% typeId: 1
%% type: 单选题
%% chapter: 直线运动
%% point: 匀变速直线运动
%% method: 无
%% score: 6
%% degree: 650
%% duplicateId: 0
%% body:
质点做直线运动的位移 $x$ 与时间 $t$ 的关系为 $x=5t+t^{2}$（各物理量均采用国际单位制单位），则该质点\xzanswer{D}
\fourchoices[answer=D]
{第 $1\Us$ 内的位移是 $5\Um$}
{前 $2\Us$ 内的平均速度是 $6\Ums$}
{任意相邻 $1\Us$ 内的位移差都是 $1\Um$}
{任意 $1\Us$ 内的速度增量都是 $2\Ums$}
%% answer: D
%% memo:
\memoanswer{
第 $1\Us$ 内的位移只需将 $t=1$ 代入即可求出 $x=6\Um$，A 错误；前 $2\Us$ 内的平均速度为 $\overline{v}=\frac{s_2}{2}=\frac{5\times2+2^{2}}{2}=7\Ums$，B 错；由题给解析式可以求得加速度为 $a=2\Umsq$，$\Delta x=aT^{2}=2\Um$，C 错；由加速度的定义可知 D 正确。
}

\item
%% number: 4
%% paperName: 2011年天津理综第 4 题 6 分
%% typeId: 1
%% type: 单选题
%% chapter: 交流电
%% point: 交流电
%% method: 无
%% score: 6
%% degree: 450
%% duplicateId: 0
%% body:
在匀强磁场中，一矩形金属线框绕与磁感线垂直的转轴匀速转动，如图 \ref{2011天津理综04a} 所示，产生的交变电动势的图象如图 \ref{2011天津理综04b} 所示，则\xzanswer{B}

\twopicture[labelA=2011天津理综04a, labelB=2011天津理综04b, numA=a, numB=b, h=4cm]{figs/04a.png}{figs/04b.png}

\fourchoices[answer=B]
{$t=0.005\Us$ 时线框的磁通量变化率为零}
{$t=0.01\Us$ 时线框平面与中性面重合}
{线框产生的交变电动势有效值为 $311\UV$}
{线框产生的交变电动势的频率为 $100\UHz$}
%% answer: B
%% memo:
\memoanswer{
由图 b 可知，该交变电动势瞬时值的表达式为 $e=311\sin100\pi t$。当 $t=0.005\Us$ 时，瞬时值 $e=311\UV$，此时磁通量变化率最大，A 错；同理当 $t=0.01\Us$ 时，$e=0$，此时线框处于中性面位置，磁通量最大，磁通量的变化率为零，B 正确；对于正弦交变电流其有效值为 $E_{max}/\sqrt{2}$，题给电动势的有效值为 $220\UV$，C 错；交变电流的频率为 $f=1/T=\omega/2\pi=50\UHz$，D 错。
}

\item
%% number: 5
%% paperName: 2011年天津理综第 5 题 6 分
%% typeId: 1
%% type: 单选题
%% chapter: 电路
%% point: 电容
%% method: 无
%% score: 6
%% degree: 450
%% duplicateId: 0
%% body:
板间距为 $d$ 的平行板电容器所带电荷量为 $Q$ 时，两极板间的电势差为 $U_{1}$，板间场强为 $E_{1}$。现将电容器所带电荷量变为 $2Q$，板间距变为 $\frac{1}{2}d$，其他条件不变，这时两极板间电势差为 $U_{2}$，板间场强为 $E_{2}$，下列说法正确的是\xzanswer{C}
\fourchoices[answer=C]
{$U_{2}=U_{1}$，$E_{2}=E_{1}$}
{$U_{2}=2U_{1}$，$E_{2}=4E_{1}$}
{$U_{2}=U_{1}$，$E_{2}=2E_{1}$}
{$U_{2}=2U_{1}$，$E_{2}=2E_{1}$}
%% answer: C
%% memo:
\memoanswer{
考查平行板电容器的相关知识。$U_1=\frac{Q}{C}=\frac{4\pi kdQ}{\varepsilon S}$，$E_1=\frac{U_1}{d}=\frac{4\pi kQ}{\varepsilon S}$，当电荷量变为 $2Q$ 时，$U_2=\frac{2Q}{C}=\frac{4\pi kdQ}{\varepsilon S}=U_1$，$E_2=\frac{U_2}{d/2}=\frac{8\pi kQ}{\varepsilon S}=2E_1$，C 选项正确。
}

\item
%% number: 6
%% paperName: 2011年天津理综第 6 题 6 分
%% typeId: 2
%% type: 多选题
%% chapter: 光学
%% point: 光的干涉
%% method: 无
%% score: 6
%% degree: 650
%% duplicateId: 0
%% body:
甲、乙两单色光分别通过同一双缝干涉装置得到各自的干涉图样，设相邻两个亮条纹的中心距离为 $\Delta x$，若 $\Delta x_{\text{甲}}>\Delta x_{\text{乙}}$，则下列说法正确的是\xzanswer{BD}
\fourchoices[answer=BD]
{甲光能发生偏振现象，则乙光不能}
{真空中甲光的波长一定大于乙光的波长}
{甲光的光子能量一定大于乙光的光子能量}
{在同一种均匀介质中甲光的传播速度大于乙光}
%% answer: BD
%% memo:
\memoanswer{
偏振现象是横波特有的现象，甲、乙都可以有偏振现象发生，A 错；由 $\Delta x_{\text{甲}}>\Delta x_{\text{乙}}$，$\Delta x=\frac{L}{d}\lambda$ 可知甲光的波长大于乙光，B 正确；光子能量取决于光子的频率，而光子频率与波长成反比，C 错；波速与波长之间同步变化，D 正确。
}

\item
%% number: 7
%% paperName: 2011年天津理综第 7 题 6 分
%% typeId: 2
%% type: 多选题
%% chapter: 机械波
%% point: 波的图象
%% method: 无
%% score: 6
%% degree: 450
%% duplicateId: 0
%% body:
位于坐标原点处的波源 A 沿 $y$ 轴做简谐运动，A 刚好完成一次全振动时，在介质中形成的简谐横波的波形如图所示，B 是沿波传播方向上介质的一个质点，则\xzanswer{ABD}

\onepicture[w=5.5cm]{figs/07.png}

\fourchoices[answer=ABD]
{波源 A 开始振动时的运动方向沿 $y$ 轴负方向}
{此后 $\frac{1}{4}$ 周期内回复力对波源 A 一直做负功}
{经半个周期时间质点 B 将向右迁移半个波长}
{在一个周期时间内 A 所受回复力的冲量为零}
%% answer: ABD
%% memo:
\memoanswer{
由 A 刚好完成一次全振动时的图线可知波源 A 的起振方向沿 $y$ 轴负方向，A 正确；经过 $\frac{1}{4}$ 周期之后质点 A 将向负向最大位移运动，回复力做负功，B 正确；质点不随波迁移，C 错误；由简谐运动的对称性可知回复力在一个周期内的冲量为零，D 正确。
}

\item
%% number: 8
%% paperName: 2011年天津理综第 8 题 6 分
%% typeId: 2
%% type: 多选题
%% chapter: 曲线运动
%% point: 人造地球卫星
%% method: 无
%% score: 6
%% degree: 650
%% duplicateId: 0
%% body:
质量为 $m$ 的探月航天器在接近月球表面的轨道上飞行，其运动视为匀速圆周运动。已知月球质量为 $M$，月球半径为 $R$，月球表面重力加速度为 $g$，引力常量为 $G$，不考虑月球自转的影响，则航天器的\xzanswer{AC}
\fourchoices[answer=AC]
{线速度 $v=\sqrt{\frac{GM}{R}}$}
{角速度 $\omega=\sqrt{Rg}$}
{运行周期 $T=2\pi\sqrt{\frac{R}{g}}$}
{向心加速度 $a=G\frac{M}{R^{2}}$}
%% answer: AC
%% memo:
\memoanswer{
万有引力提供卫星做圆周运动的向心力，代入相关公式即可。
}

\item
%% number: 9
%% paperName: 2011年天津理综第 9 题 8 分
%% typeId: 4
%% type: 实验题
%% chapter: 电路
%% point: 伏安法测电阻
%% method: 无
%% score: 8
%% degree: 450
%% duplicateId: 0
%% body:
（18 分）

\begin{enumerate}
	\item
	某同学用测力计研究在竖直方向运行的电梯运动状态。他在地面上用测力计测量砝码的重力，示数为 $G$。他在电梯中用测力计仍测量同一砝码的重力，发现测力计的示数小于 $G$，由此判断此时电梯的运动状态可能是\underline{\hspace{2cm}}。

	\item
	用螺旋测微器测量某金属丝直径的结果如图所示。该金属丝的直径是\underline{\hspace{1.6cm}}$\Umm$。

	\onepicture[align=right, w=4.5cm]{figs/09a.png}
	\item
	某同学用大头针、三角板、量角器等器材测半圆形玻璃砖的折射率。开始玻璃砖的位置如图中实线所示，使大头针 P$_{1}$、P$_{2}$ 与圆心 O 在同一直线上，该直线垂直于玻璃砖的直径边，然后使玻璃砖绕圆心 O 缓慢转动，同时在玻璃砖直径边一侧观察 P$_{1}$、P$_{2}$ 的像，且 P$_{2}$ 的像挡住 P$_{1}$ 的像。如此观察，当玻璃砖转到图中虚线位置时，上述现象恰好消失。此时只需测量出\underline{\hspace{2.4cm}}，即可计算出玻璃砖的折射率。请用你的测量量表示出折射率\underline{\hspace{1.6cm}}。

	\onepicture[align=right, w=4.5cm]{figs/09b.png}
	\item
	某同学测量阻值约为 $25\UkO$ 的电阻 $R_{x}$，现备有下列器材：

	A．电流表（量程 $100\UuA$，内阻约为 $2\UkO$）；

	B．电流表（量程 $500\UuA$，内阻约为 $300\UO$）；

	C．电压表（量程 $15\UV$，内阻约为 $100\UkO$）；

	D．电流表（量程 $50\UV$，内阻约为 $500\UkO$）；

	E．直流电源（$20\UV$，允许最大电流 $1\UA$）；

	F．滑动变阻器（最大阻值 $1\UkO$，额定功率 $1\UW$）；

	G．电键和导线若干。

	电流表应选\underline{\hspace{1.6cm}}，电压表应选\underline{\hspace{1.6cm}}。（填字母代号）

	该同学正确选择仪器后连接了以下电路，为保证实验顺利进行，并使测量误差尽量减小，实验前请你检查该电路，指出电路在接线上存在的问题：

	\onepicture[align=right, w=6.5cm]{figs/09c.png}

	\begin{steps}
		\item
		\underline{\hspace{4cm}}；
		\item
		\underline{\hspace{4cm}}。
	\end{steps}
\end{enumerate}
%% answer: 
\jdanswer{
\begin{enumerate}
	\item
	减速上升或加速下降
	
	\item
	$1.706\Umm$
	
	\item
	玻璃砖直径边绕 O 点转过的角度，$n=\frac{1}{\sin\theta}$
	
	\item
	B、C；①电流表应采用内接的方法；②滑动变阻器应采用分压器方式的接法。
\end{enumerate}
}
%% memo:
\memoanswer{
（1）物体处于失重状态，加速度方向向下，故而可能是减速上升或加速下降。

（2）注意副尺一定要有估读。读数为 $1.5+20.6\times0.01\Umm=1.706\Umm$。因为个人情况不同，估读不一定一致，本题读数 $1.704\sim1.708$ 都算正确。

（3）由题意可知，当玻璃砖转过某一角度 $\theta$ 时，刚好发生全反射，在直径边一侧观察不到 P$_{1}$、P$_{2}$ 的像，作出如图所示的光路图可知，当转过角度 $\theta$ 时有 $n=\frac{1}{\sin\theta}$。

\onepicture[w=4.5cm]{figs/09_an.png}

（4）电学实验选择仪器的一般步骤如下：①根据量程选择电流表和电压表，不能超过表的量程，不能量程太大导致表的读数偏小；②根据题中关键语句，如精确测量，从零开始连续可调等等选择分压电路亦或是限流电路；分压电路滑动变阻器选择小阻值，限流电路滑动变阻器选择大阻值；③选择电流表的内外接法，一般的原则是“大内偏大，小外偏小”；也可以根据 $R_{V}/R_{x}$ 与 $R_{x}/R_{A}$ 之间的关系来判断，当 $R_{V}/R_{x}>R_{x}/R_{A}$ 时，采用电流表的外接法，反之选择电流表内接法。

本题中，待测电阻 $R_x$ 的阻值约为 $25\UkO$，直流电源电动势为 $20\UV$，经粗略计算电路中最大的电流约为 $I_{max}=\frac{E}{R}=\frac{20}{25\times10^{3}}=800\UuA$，所以电流表选择 B；虽然电压表 C 的量程不足，但是相比起来电压表 D 的量程超过太多，读数偏小，所以电压表选择 C 表。

根据本题解析的第②③两条原则可知电路图的问题为：①电流表应采用内接的方法；②滑动变阻器应采用分压器方式的接法。
}

\item
%% number: 10
%% paperName: 2011年天津理综第 10 题 16 分
%% typeId: 6
%% type: 计算题
%% chapter: 曲线运动
%% point: 平抛运动
%% method: 无
%% score: 16
%% degree: 450
%% duplicateId: 0
%% body:
如图所示，圆管构成的半圆形竖直轨道固定在水平地面上，轨道半径为 $R$，MN 为直径且与水平面垂直，直径略小于圆管内径的小球 A 以某一初速度冲进轨道，到达半圆轨道最高点 M 时与静止于该处的质量与 A 相同的小球 B 发生碰撞，碰后两球粘在一起飞出轨道，落地点距 N 为 $2R$。重力加速度为 $g$，忽略圆管内径，空气阻力及各处摩擦均不计，求：

\begin{enumerate}
	\item
	粘合后的两球从飞出轨道到落地的时间 $t$；
	\item
	小球 A 冲进轨道时速度 $v$ 的大小。
\end{enumerate}

\onepicture[align=right, w=5.5cm]{figs/10.png}
%% answer: 
\jdanswer{
\begin{enumerate}
	\item
	$2\sqrt{\frac{R}{g}}$
	
	\item
	$2\sqrt{2gR}$
\end{enumerate}
}
%% memo:
\memoanswer{
（1）粘合后的两球飞出轨道后做平抛运动，竖直方向分运动为自由落体运动，有 $2R=\frac{1}{2}gt^{2}$，解得 $t=2\sqrt{\frac{R}{g}}$。

（2）设球 A 的质量为 $m$，碰撞前速度大小为 $v_1$，把球 A 冲进轨道最低点时的重力势能定为 $0$，由机械能守恒定律知 $\frac{1}{2}mv^{2}=\frac{1}{2}mv_1^{2}+2mgR$，设碰撞后粘合在一起的两球速度大小为 $v_2$，由动量守恒定律知 $mv_1=2mv_2$，飞出轨道后做平抛运动，水平方向分运动为匀速直线运动，有 $2R=v_2t$，综合得 $v=2\sqrt{2gR}$。
}

\item
%% number: 11
%% paperName: 2011年天津理综第 11 题 18 分
%% typeId: 6
%% type: 计算题
%% chapter: 电磁感应
%% point: 电磁感应能量分析
%% method: 无
%% score: 18
%% degree: 450
%% duplicateId: 0
%% body:
如图所示，两根足够长的光滑金属导轨 MN、PQ 间距为 $l=0.5\Um$，其电阻不计，两导轨及其构成的平面均与水平面成 $30^\circ$ 角。完全相同的两金属棒 ab、cd 分别垂直导轨放置，每棒两端都与导轨始终有良好接触，已知两棒的质量均为 $0.02\Ukg$，电阻均为 $R=0.1\UO$，整个装置处在垂直于导轨平面向上的匀强磁场中，磁感应强度为 $B=0.2\UT$，棒 ab 在平行于导轨向上的力 $F$ 作用下，沿导轨向上匀速运动，而棒 cd 恰好能保持静止。取 $g=10\Umsq$，问：

\begin{enumerate}
	\item
	通过 cd 棒的电流 $I$ 是多少，方向如何？
	\item
	棒 ab 受到的力 $F$ 多大？
	\item
	棒 cd 每产生 $Q=0.1\UJ$ 的热量，力 $F$ 做的功 $W$ 是多少？
\end{enumerate}

\onepicture[align=right, w=6.5cm]{figs/11.png}
%% answer: 
\jdanswer{
\begin{enumerate}
	\item
	$1\UA$，方向由 d 到 c
	
	\item
	$0.2\UN$
	
	\item
	$0.4\UJ$
\end{enumerate}
}
%% memo:
\memoanswer{
（1）棒 cd 受到的安培力 $F_{cd}=IlB$，棒 cd 在共点力作用下平衡，则 $F_{cd}=mg\sin30^\circ$，代入数据解得 $I=1\UA$，方向由右手定则可知由 d 到 c。

（2）棒 ab 与棒 cd 受到的安培力大小相等 $F_{ab}=F_{cd}$，对棒 ab 由共点力平衡有 $F=mg\sin30^\circ+IlB$，代入数据解得 $F=0.2\UN$。

（3）设在时间 $t$ 内棒 cd 产生 $Q=0.1\UJ$ 热量，由焦耳定律可知 $Q=I^{2}Rt$，设 ab 棒匀速运动的速度大小为 $v$，则产生的感应电动势 $E=Blv$，由闭合电路欧姆定律知 $I=\frac{E}{2R}$，由运动学公式知，在时间 $t$ 内，棒 ab 沿导轨的位移 $x=vt$，力 $F$ 做的功 $W=Fx$，综合上述各式，代入数据解得 $W=0.4\UJ$。
}

\item
%% number: 12
%% paperName: 2011年天津理综第 12 题 20 分
%% typeId: 6
%% type: 计算题
%% chapter: 磁场
%% point: 带电粒子在磁场中的运动
%% method: 无
%% score: 20
%% degree: 250
%% duplicateId: 0
%% body:
回旋加速器在核科学、核技术、核医学等高新技术领域得到了广泛应用，有力地推动了现代科学技术的发展。

\begin{enumerate}
	\item
	当今医学成像诊断设备 PET/CT 堪称“现代医学高科技之冠”，它在医疗诊断中，常利用能放射电子的同位素碳 11 为示踪原子，碳 11 是由小型回旋加速器输出的高速质子轰击氮 14 获得，同时还产生另一粒子，试写出核反应方程。若碳 11 的半衰期 $\tau$ 为 $20\Umin$，经 $2.0\Uh$ 剩余碳 11 的质量占原来的百分之几？（结果取 2 位有效数字）
	\item
	回旋加速器的原理如图，D$_{1}$ 和 D$_{2}$ 是两个中空的半径为 $R$ 的半圆金属盒，它们接在电压一定、频率为 $f$ 的交流电源上，位于 D$_{1}$ 圆心处的质子源 A 能不断产生质子（初速度可以忽略，重力不计），它们在两盒之间被电场加速，D$_{1}$、D$_{2}$ 置于与盒面垂直的磁感应强度为 $B$ 的匀强磁场中。若质子束从回旋加速器输出时的平均功率为 $P$，求输出时质子束的等效电流 $I$ 与 $P$、$B$、$R$、$f$ 的关系式（忽略质子在电场中运动的时间，其最大速度远小于光速）
	\item
	试推理说明：质子在回旋加速器中运动时，随轨道半径 $r$ 的增大，同一盒中相邻轨道的半径之差 $\Delta r$ 是增大、减小还是不变？
\end{enumerate}

\onepicture[align=right, w=6cm]{figs/12.png}
%% answer: 
\jdanswer{
\begin{enumerate}
	\item
	$1.6\%$
	
	\item
	$I=\frac{P}{\pi BR^{2}f}$
	
	\item
	随轨道半径 $r$ 的增大，同一盒中相邻轨道的半径之差 $\Delta r$ 减小
\end{enumerate}
}
%% memo:
\memoanswer{
（1）核反应方程为 \ce{^{14}_{7}N + ^{1}_{1}H -> ^{11}_{6}C + ^{4}_{2}He}，设碳 11 原有质量为 $m_0$，经过 $t=2.0\Uh$ 剩余的质量为 $m_t$，根据半衰期定义，有：$m_t/m_0=(1/2)^{t/\tau}=(1/2)^{120/20}=1.6\%$。

（2）设质子质量为 $m$，电荷量为 $q$，质子离开加速器时速度大小为 $v$，由牛顿第二定律知：$qvB=m\frac{v^{2}}{R}$，质子运动的回旋周期为：$T=\frac{2\pi R}{v}=\frac{2\pi m}{qB}$，由回旋加速器工作原理可知，交变电源的频率与质子回旋频率相同，由周期 $T$ 与频率 $f$ 的关系可得：$f=\frac{1}{T}$，设在 $t$ 时间内离开加速器的质子数为 $N$，则质子束从回旋加速器输出时的平均功率 $P=\frac{Nmv^{2}}{2t}$，输出时质子束的等效电流为：$I=\frac{Nq}{t}$，由上述各式得 $I=\frac{P}{\pi BR^{2}f}$。

（3）方法一：设 $k(k\in N^{*})$ 为同一盒子中质子运动轨道半径的序数，相邻的轨道半径分别为 $r_k$、$r_{k+1}(r_{k+1}>r_k)$，$\Delta r_k=r_{k+1}-r_k$，在相应轨道上质子对应的速度大小分别为 $v_k$、$v_{k+1}$，D$_{1}$、D$_{2}$ 之间的电压为 $U$，由动能定理知 $2qU=\frac{1}{2}mv_{k+1}^{2}-\frac{1}{2}mv_k^{2}$，由洛伦兹力充当质子做圆周运动的向心力，知 $r_k=\frac{mv_k}{qB}$，则 $2qU=\frac{q^{2}B^{2}}{2m}(r_{k+1}^{2}-r_k^{2})$，整理得 $\Delta r_k=\frac{4mU}{qB^{2}(r_{k+1}+r_k)}$，因 $U$、$q$、$m$、$B$ 均为定值，令 $C=\frac{4mU}{qB^{2}}$，由上式得 $\Delta r_k=\frac{C}{r_{k+1}+r_k}$，相邻轨道半径 $r_{k+1}$、$r_{k+2}$ 之差 $\Delta r_{k+1}=r_{k+2}-r_{k+1}$，同理 $\Delta r_{k+1}=\frac{C}{r_{k+1}+r_{k+2}}$，因为 $r_{k+2}>r_k$，比较 $\Delta r_k$、$\Delta r_{k+1}$ 得 $\Delta r_{k+1}<\Delta r_k$，说明随轨道半径 $r$ 的增大，同一盒中相邻轨道的半径之差 $\Delta r$ 减小。

方法二：设 $k(k\in N^{*})$ 为同一盒子中质子运动轨道半径的序数，相邻的轨道半径分别为 $r_k$、$r_{k+1}(r_{k+1}>r_k)$，$\Delta r_k=r_{k+1}-r_k$，在相应轨道上质子对应的速度大小分别为 $v_k$、$v_{k+1}$，D$_{1}$、D$_{2}$ 之间的电压为 $U$，由洛伦兹力充当质子做圆周运动的向心力，知 $r_k=\frac{mv_k}{qB}$，故 $r_k/r_{k+1}=v_k/v_{k+1}$，由动能定理知，质子每加速一次，其动能增量 $\Delta E_k=qU$，以质子在 D$_{2}$ 盒中运动为例，第 $k$ 次进入 D$_{2}$ 时，被电场加速 $(2k-1)$ 次，速度大小为 $v_k=\sqrt{\frac{(2k-1)2qU}{m}}$，同理，质子第 $k+1$ 次进入 D$_{2}$ 时，速度大小为 $v_{k+1}=\sqrt{\frac{(2k+1)2qU}{m}}$，综合上述各式可得 $r_k/r_{k+1}=v_k/v_{k+1}=\sqrt{\frac{2k-1}{2k+1}}$，整理得 $r_k^{2}/r_{k+1}^{2}=\frac{2k-1}{2k+1}$，$\Delta r_k=\frac{2r_{k+1}^{2}}{(2k+1)(r_{k+1}+r_k)}$，同理，对于相邻轨道半径 $r_{k+1}$、$r_{k+2}$，$\Delta r_{k+1}=r_{k+2}-r_{k+1}$，整理后有 $\Delta r_{k+1}=\frac{2r_{k+1}^{2}}{(2k+1)(r_{k+1}+r_{k+2})}$，因为 $r_{k+2}>r_k$，比较 $\Delta r_k$、$\Delta r_{k+1}$ 得 $\Delta r_{k+1}<\Delta r_k$，说明随轨道半径 $r$ 的增大，同一盒中相邻轨道的半径之差 $\Delta r$ 减小。
}

\end{enumerate}

\end{document}
